% Generated draft: requires mathematical and source review.
\documentclass[11pt]{article}
\usepackage[margin=25mm]{geometry}
\usepackage{fontspec}
\setmainfont{Latin Modern Roman}
\usepackage{amsmath,amssymb,mathtools,unicode-math}
\setmathfont{Latin Modern Math}
\usepackage{enumitem,needspace,xurl,hyperref}
\setlength{\parindent}{0pt}
\setlength{\parskip}{5pt}
\setlength{\emergencystretch}{3em}
\newcommand{\N}{\mathbb{N}}
\newcommand{\Z}{\mathbb{Z}}
\newcommand{\Q}{\mathbb{Q}}
\newcommand{\R}{\mathbb{R}}
\newcommand{\C}{\mathbb{C}}
\newcommand{\F}{\mathbb{F}}
\newcommand{\A}{\mathbb{A}}
\newcommand{\PP}{\mathbb{P}}
\newcommand{\E}{\mathbb{E}}
\newcommand{\eps}{\varepsilon}
\newcommand{\dd}{\,\mathrm{d}}
\begin{document}
\section*{UnsolvedMath ID: 701090}
\textbf{Source problem code:} catalogue entry
\subsection{Definitions and mathematical statement}
Let \(G\) be a finite group, and let \(\operatorname{Irr}(G)\) denote the set of characters of its irreducible complex representations. An element \(x\in G\) is called \emph{nonvanishing} if
\[
\chi(x)\ne 0\qquad\text{for every }\chi\in\operatorname{Irr}(G).
\]
Write \(\mathrm{NV}(G)\) for the set of these elements, and write \(F(G)\) for the Fitting subgroup of \(G\), the largest nilpotent normal subgroup of \(G\). The order of \(x\) is denoted by \(|x|\).

The source recalls Burnside's theorem that every \(\chi\in\operatorname{Irr}(G)\) with \(\chi(1)>1\) vanishes at some element of \(G\). Thus the irreducible characters that are nonzero everywhere are precisely the linear characters. The questions below concern the dual condition on elements.

\begin{enumerate}
\item[(a)] Is it true that for every finite solvable group \(G\) and every \(x\in G\),
\[
\bigl(x\in\mathrm{NV}(G)\ \text{and}\ |x|\equiv 0\pmod 2\bigr)
\quad\Longrightarrow\quad x\in F(G)?
\]
The source attributes the corresponding assertion for elements of odd order to I. M. Isaacs, G. Navarro, and T. R. Wolf (1999). Consequently, the question is equivalently whether \(\mathrm{NV}(G)\subseteq F(G)\) for every finite solvable group \(G\).
\item[(b)] Determine, up to isomorphism, all finite nonabelian simple groups \(S\) for which
\[
\exists x\in S\setminus\{1\}\ \forall\chi\in\operatorname{Irr}(S),\qquad \chi(x)\ne 0.
\]
The source gives \(\mathbb{A}_7\), the alternating group on seven letters, as an example having such an element.
\end{enumerate}

The source also attributes to M. Miyamoto (2008) the assertion that every nontrivial abelian normal subgroup of a finite group contains a nonvanishing element. As supplied, this assertion does not require that element to be nonidentity; the identity is always nonvanishing since \(\chi(1)\) is a positive integer. No stronger version of this ancillary assertion is assumed here. The registry labels the problem open; that label and the supplied background assertions have not been independently verified.
\subsection{Short English statement}
For a finite group, an element is nonvanishing if every irreducible complex character is nonzero on it. Part (a) asks whether every nonvanishing element of even order in a finite solvable group lies in its Fitting subgroup, extending the supplied odd-order result. Part (b) asks which finite nonabelian simple groups contain a nonidentity nonvanishing element. The source lists \(\mathbb{A}_7\) as an example. This is a statement-only draft requiring human mathematical review.
\subsection{Sources}
\begin{itemize}
\item Local UnsolvedMath registry: supplied complete source record, permanent ID \(701090\), exact source problem code \texttt{catalogue entry}, title \emph{Kourovka Notebook Problem 15.2}. The supplied registry status is \texttt{open}; it is not an independently verified literature assessment.
\item \emph{The Kourovka Notebook (Unsolved Problems in Group Theory)}, Problem 15.2, proposed by I. M. Isaacs in the 15th issue (2002). The supplied record identifies the 21st issue edition (2026), \texttt{arXiv:1401.0300v48}, and describes the collection as edited by E. I. Khukhro and V. D. Mazurov. These bibliographic details are supplied source information, not independently verified.
\item I. M. Isaacs, G. Navarro, and T. R. Wolf, \emph{J. Algebra} \textbf{222}, no. 2 (1999), 413--423. The record cites this reference for the odd-order assertion in part (a); no article title is supplied.
\item W. Burnside: attribution in the supplied statement for the theorem that every nonlinear irreducible complex character of a finite group has a zero. No separate bibliographic citation is supplied.
\item M. Miyamoto (2008): attribution in the supplied statement for the assertion concerning nontrivial abelian normal subgroups. No publication citation is supplied, and the supplied wording does not require a nonidentity element.
\end{itemize}
\end{document}
